\chapter{Rozwinięcie}
\label{ch:rozwiniecie}

W~rozdziale opisano przedmiot i~metodykę badań. Przedstawiono cel badań,
przebieg eksperymentu, sposób przygotowania danych uczących i~testowych,
porównywane warianty trasowania oraz miary, według których je oceniano.

% TODO: część dotycząca RL_PCB (rozmieszczanie elementów).

\section{Cel badań}
\label{sec:cel-badan}

Celem badań jest ocena, czy rozwiązanie Unet-Astar~\cite{Yin2023UnetAstar}
wytrenowane wyłącznie na syntetycznych, losowo generowanych płytkach może
skutecznie wspierać trasowanie rzeczywistych projektów PCB. Sformułowano
następujące pytania badawcze:
\begin{enumerate}
  \item Czy predykcja sieci U-Net zmniejsza liczbę węzłów rozwijanych przez
    algorytm A* na płytkach rzeczywistych, a~więc czy sieć uogólnia się poza
    rozkład danych uczących?
  \item Jak jakość trasowania uzyskanego algorytmem Unet-Astar (odsetek
    połączonych sieci, długość ścieżek, liczba przelotek, liczba naruszeń reguł
    projektowych) wypada na tle autoroutera wbudowanego w~EasyEDA?
  \item Czy różnice w~geometrii ścieżek przekładają się na właściwości istotne
    z~punktu widzenia kompatybilności elektromagnetycznej?
\end{enumerate}

\section{Przebieg badań}
\label{sec:przebieg-badan}

Algorytm Unet-Astar nie wyznacza ścieżek samodzielnie --- sieć neuronowa
jedynie zmienia kolejność rozwijania węzłów przez A*
(wzór~\ref{eq:unet-astar-f}). Wpływ sieci na jakość trasowania jest więc
pośredni, a~jej główną zaletą ma być skrócenie wyszukiwania. Z~tego powodu
porównanie wyłącznie z~autorouterem EasyEDA nie pozwoliłoby rozdzielić wpływu
sieci od wpływu samego algorytmu A*. W~badaniach porównano zatem trzy warianty
trasowania:
\begin{description}
  \item[EasyEDA] --- wbudowany autorouter EasyEDA (wariant odniesienia);
  \item[A*] --- algorytm wielowarstwowego A* z~repozytorium Unet-Astar
    uruchomiony bez sieci neuronowej;
  \item[Unet-A*] --- ten sam algorytm wspomagany predykcją sieci U-Net.
\end{description}

Badania przebiegały w~następujących krokach:
\begin{enumerate}
  \item Wygenerowanie zbioru uczącego z~losowych płytek
    (sekcja~\ref{sec:zbior-uczacy}).
  \item Trening sieci U-Net (sekcja~\ref{sec:trening}).
  \item Przygotowanie płytek testowych w~EasyEDA: projektów własnych oraz
    zaimportowanych projektów gotowych (sekcja~\ref{sec:zbior-testowy}).
  \item Trasowanie każdej płytki testowej autorouterem EasyEDA i~eksport
    pliku autoroutera w~formacie \emph{Specctra DSN}.
  \item Trasowanie tej samej płytki na podstawie pliku DSN w~wariantach A*
    i~Unet-A*.
  \item Import wyniku (plik \emph{Specctra SES}) z~powrotem do EasyEDA,
    weryfikacja DRC i~odczyt statystyk projektu.
  \item Zestawienie miar (sekcja~\ref{sec:miary}) dla wszystkich płytek
    testowych i~wariantów.
\end{enumerate}

Wszystkie warianty otrzymują identyczne dane wejściowe: obrys płytki, rozmieszczenie
elementów, sieci oraz reguły projektowe zapisane w~tym samym pliku DSN.

\section{Środowisko badawcze}
\label{sec:srodowisko}

\begin{table}[ht]
\fcmtcaption{Środowisko badawcze. Źródło: opracowanie własne.}\label{tab:srodowisko}
\centering\footnotesize%
\begin{tabular}{p{0.3\textwidth} p{0.6\textwidth}}
\toprule
Element & Wersja / parametry \\
\midrule
System operacyjny & Windows 11 Home \\
\midrule
Procesor & Intel Core i5-8400, 2,8~GHz, 6~rdzeni / 6~wątków \\
\midrule
Pamięć RAM & 24~GB, 2133~MT/s \\
\midrule
Karta graficzna & NVIDIA GeForce GTX 1060, 3~GB \\
\midrule
Python & 3.8.10 \\
\midrule
PyTorch & 2.4.1 (CUDA 12.1) \\
\midrule
EasyEDA & Standard, wersja 6.5.61 \\
\midrule
Unet-Astar & repozytorium~\cite{Firesuiry_UnetAstar_GitHub} wraz
  z~modyfikacjami opisanymi w~sekcji~\ref{sec:modyfikacje} \\
\bottomrule
\end{tabular}
\end{table}

\section{Format danych i~współpraca z~EasyEDA}
\label{sec:format-danych}

Wymiana danych z~EasyEDA odbywa się za pośrednictwem plików autoroutera
w~formacie \emph{Specctra}. Plik DSN zawiera obrys płytki, definicje warstw,
reguły projektowe (odstęp izolacyjny, szerokość ścieżki, przelotki), kształty
padów (\emph{padstack}), położenia pinów oraz listę sieci. EasyEDA zapisuje
płytkę jako jeden komponent, którego piny mają współrzędne bezwzględne.
Obroty elementów są uwzględnione bezpośrednio w~kształtach padów.
Współrzędne zapisane są w~jednostkach EasyEDA równych 10~mil.

Program Unet-Astar przelicza współrzędne na mile i~na potrzeby trasowania
przenosi je na siatkę o~boku 8~mil (ok.~0{,}2~mm). Reguły projektowe są przy
tym zaokrąglane w~górę do całkowitej liczby oczek siatki.

Oryginalny program Unet-Astar nie zapisuje wyniku w~formacie, który można
wczytać z~powrotem do EasyEDA. Wymiana danych przebiega więc w~dwóch
kierunkach: z~EasyEDA eksportowany jest plik DSN, będący danymi wejściowymi
programu Unet-Astar, a~wynik trasowania wraca do EasyEDA jako plik sesji
\emph{Specctra SES} importowany do projektu.
Plik SES generowany jest przez eksporter przygotowany na potrzeby badań.
Eksporter odwzorowuje kolejne punkty ścieżki z~siatki
z~powrotem na współrzędne płytki, zmianę warstwy zapisuje jako przelotkę,
a~ścieżki podsieci powstałych z~podziału sieci wielopinowych łączy pod
wspólną nazwą sieci z~pliku DSN. Szerokość ścieżki i~średnicę przelotki
przyjmuje się z~reguł zapisanych w~pliku DSN, bez zaokrąglenia do siatki.
Ponieważ siatka o~boku 8~mil nie pokrywa się ze środkami padów, końce
ścieżek dociągane są krótkim odcinkiem do środka padu. Odcinek ten
może naruszać reguły odstępów, dlatego zaimportowana płytka jest
sprawdzana narzędziem DRC programu EasyEDA.

\section{Zbiór uczący}
\label{sec:zbior-uczacy}

Sieć trenowano wyłącznie na płytkach syntetycznych, generowanych losowo zgodnie
z~procedurą z~repozytorium Unet-Astar. Parametry generatora zestawiono
w~tabeli~\ref{tab:generator}.

\begin{table}[ht]
\fcmtcaption{Parametry generatora losowych płytek. Źródło: opracowanie własne
na podstawie kodu źródłowego~\cite{Firesuiry_UnetAstar_GitHub}.}\label{tab:generator}
\centering\footnotesize%
\begin{tabular}{p{0.45\textwidth} p{0.45\textwidth}}
\toprule
Parametr & Wartość \\
\midrule
Liczba płytek & 1250 \\
\midrule
Wymiary płytki & losowo 1000--2000~mil (25{,}4--50{,}8~mm) w~każdej osi \\
\midrule
Liczba warstw & 2 \\
\midrule
Gęstość pinów & 0{,}15 pinu na pole $100 \times 100$~mil \\
\midrule
Gęstość przeszkód & 0{,}1 przeszkody na pole $300 \times 300$~mil \\
\midrule
Kształty padów & koło o~średnicy 80~mil lub prostokąt $56 \times 48$~mil \\
\midrule
Przeszkody & koło o~średnicy 240~mil, bez przypisanej sieci \\
\midrule
Minimalne odległości & pin--pin 100~mil, pin--przeszkoda 300~mil,
  przeszkoda--przeszkoda 500~mil \\
  \midrule
Sieci & wyłącznie dwupinowe, piny przydzielane losowo \\
\midrule
Siatka trasowania & 8~mil \\
\bottomrule
\end{tabular}
\end{table}

Górną granicę wymiarów płytki dobrano tak, aby płytka mieściła się w~wejściu
sieci o~rozmiarze $256 \times 256$ oczek siatki (2048~mil). W~oryginalnym
repozytorium generowano płytki do 5000~mil, jednak próbki większe od wejścia
sieci są odrzucane przez moduł wczytujący dane.

Każda wygenerowana płytka jest najpierw w~całości trasowana algorytmem A*.
Sieci są prowadzone kolejno, od najkrótszej. Następnie dla każdego połączenia
(z~pominięciem 30\% najkrótszych) tworzona jest próbka ucząca:
\begin{description}
  \item[Wejście] --- obraz o~trzech kanałach. Kanał pierwszy zawiera punkt
    początkowy i~końcowy połączenia rozmyte filtrem Gaussa o~odchyleniu
    standardowym równym 30\% odległości między nimi. Kanały drugi i~trzeci
    zawierają mapy przeszkód na obu warstwach: pady innych sieci oraz
    pozostałe, już poprowadzone ścieżki.
  \item[Etykieta] --- dwukanałowa maska obszaru rekomendowanego,
    wyznaczona dwukierunkowym algorytmem Lee (sekcja~\ref{sec:lee}).
\end{description}

Algorytm Lee służy wyłącznie do przygotowania etykiet. Wyznacza obszar
rekomendowany dokładnie, ale przeszukuje przy tym całą dostępną część
płytki, więc nie przyspieszyłby trasowania. Sieć uczy się odtwarzać jego
wynik: kosztowne obliczenia wykonywane są jednokrotnie, podczas generowania
zbioru uczącego, a~w~trakcie trasowania zastępuje je pojedyncza predykcja
sieci. Jakość obszarów wyznaczonych przez algorytm Lee stanowi zatem górną
granicę tego, czego sieć może się nauczyć.

Implementacja w~repozytorium różni się od opisu w~artykule
(wzór~\ref{eq:lee-obszar}). Do obszaru rekomendowanego należą oczka
spełniające warunek $d_s(n) + d_t(n) \le L^* + 0{,}1\,d_M(s,t) + 1$, gdzie
$d_M$ oznacza odległość w~metryce miejskiej zamiast euklidesowej, a~stały
składnik $+1$ poszerza obszar o~jedno oczko zapasu. W~badaniach
wykorzystano etykiety wyznaczone zgodnie z~implementacją.

Uzyskano 13\,830 próbek. Zbiór podzielono losowo w~proporcji
90:10 na część uczącą i~walidacyjną.

\section{Model i~trening}
\label{sec:trening}

Repozytorium zawiera trzy architektury: klasyczny U-Net (31~mln parametrów),
U-Net++ (47~mln) oraz głęboką sieć \emph{ResNetUNet} zaproponowaną przez
autorów rozwiązania (170~mln parametrów, 8~poziomów redukcji rozdzielczości).
Próba treningu modelu \emph{ResNetUNet} na dostępnej karcie graficznej
z~3~GB pamięci wykazała, że już przy dwóch próbkach w~paczce zapotrzebowanie
na pamięć przekracza jej pojemność. Trening trwał wówczas ponad 4~s na
paczkę. Z~tego powodu w~badaniach wykorzystano klasyczny U-Net, który przy
czterech próbkach w~paczce zajmuje ok.~2~GB pamięci.
% TODO: Google Colab -- dopisać.

\begin{table}[ht]
\fcmtcaption{Parametry treningu. Źródło: opracowanie własne.}\label{tab:trening}
\centering\footnotesize%
\begin{tabular}{p{0.45\textwidth} p{0.45\textwidth}}
\toprule
Parametr & Wartość \\
\midrule
Architektura & U-Net, wejście $3 \times 256 \times 256$, wyjście
  $2 \times 256 \times 256$ \\
  \midrule
Funkcja straty & $0{,}5 \cdot \mathrm{BCE} + 0{,}5 \cdot \mathrm{Dice}$ \\
\midrule
Optymalizator & Adam, współczynnik uczenia $10^{-4}$ \\
\midrule
Harmonogram & mnożenie współczynnika uczenia przez 0{,}75 po każdej epoce \\
\midrule
Liczba epok & 15 \\
\midrule
Rozmiar paczki & 4 \\
\midrule
Wybór modelu & najmniejsza strata na zbiorze walidacyjnym \\
\bottomrule
\end{tabular}
\end{table}

W~trakcie trasowania sieć wywoływana jest dla każdego połączenia
z~wyjątkiem 30\% najkrótszych. Dla płytek większych od wejścia sieci
z~mapy przeszkód wycinany jest fragment $256 \times 256$ oczek
wyśrodkowany między punktami połączenia. Połączenia, których rozpiętość
w~którejkolwiek osi przekracza 256 oczek, prowadzone są bez udziału sieci.
Wynik sieci $s(n)$ modyfikuje funkcję oceny węzła zgodnie ze
wzorem~\ref{eq:unet-astar-f}, z~parametrami $a = 800$ i~$b = 0$,
przyjętymi za eksperymentami autorów.

\section{Zbiór testowy}
\label{sec:zbior-testowy}

Zbiór testowy obejmuje wyłącznie płytki rzeczywiste, niewykorzystane
w~treningu:
\begin{itemize}
  \item płytki zaprojektowane samodzielnie w~EasyEDA,
  \item gotowe projekty udostępnione publicznie w~serwisie EasyEDA ---
    projekt przykładowy \emph{Project-Sample} z~sekcji
    \emph{Modules}~\cite{EasyEDA_ProjectSample}.
\end{itemize}

Z~zaimportowanych projektów przed trasowaniem usuwane są ścieżki
i~przelotki poprowadzone przez ich autorów, tak aby wszystkie warianty
trasowały płytkę od początku. Ścieżki zapisane w~pliku DSN (sekcja
\emph{wiring}) są przez program Unet-Astar pomijane.
Ze względu na ograniczenia rozwiązania płytki testowe muszą spełniać
następujące warunki:
\begin{itemize}
  \item dwie warstwy miedzi --- sieć przyjmuje mapy przeszkód dokładnie
    dwóch warstw;
  \item obrys prostokątny --- program uwzględnia wyłącznie prostokąt
    opisany na obrysie płytki;
  \item brak obszarów zabronionych i~stref miedzi --- nie są one
    odczytywane z~pliku DSN.
\end{itemize}
Płytki większe niż $2048 \times 2048$~mil mogą być trasowane, ale miary
jakości predykcji (sekcja~\ref{sec:miary}) wyznaczane są wyłącznie dla
płytek mieszczących się w~wejściu sieci.

Rozkład danych testowych istotnie różni się od rozkładu danych uczących.
Płytki rzeczywiste zawierają sieci wielopinowe, które program dzieli na
połączenia dwupinowe. Mają też gęsto rozmieszczone pady elementów SMD
o~różnych kształtach i~nierównomierne zagęszczenie połączeń. Badanie
ma zatem charakter oceny zdolności uogólniania modelu.

\section{Miary oceny}
\label{sec:miary}

Miary zestawiono w~tabeli~\ref{tab:miary}. Miary jakości trasowania
odczytywane są w~EasyEDA po zaimportowaniu wyniku, dzięki czemu wszystkie
warianty oceniane są tym samym narzędziem. Miary wydajności wyszukiwania
dotyczą wyłącznie wariantów A* i~Unet-A*.

\begin{table}[ht]
\fcmtcaption{Miary oceny. Źródło: opracowanie własne.}\label{tab:miary}
\centering\footnotesize%
\begin{tabular}{p{0.28\textwidth} p{0.42\textwidth} p{0.2\textwidth}}
\toprule
Miara & Opis & Warianty \\
\midrule
\multicolumn{3}{l}{\emph{Jakość trasowania}} \\
\midrule
Odsetek połączeń & udział poprowadzonych połączeń w~liczbie wszystkich
  połączeń & wszystkie \\
  \midrule
Długość ścieżek & łączna długość ścieżek na obu warstwach & wszystkie \\
\midrule
Liczba przelotek & --- & wszystkie \\
\midrule
Naruszenia DRC & liczba błędów zgłoszonych przez DRC EasyEDA, w~podziale
  na typy & wszystkie \\
  \midrule
Czas trasowania & czas od wczytania płytki do zakończenia trasowania
  & wszystkie \\
\midrule
\multicolumn{3}{l}{\emph{Wydajność wyszukiwania}} \\
\midrule
Rozwinięte węzły & liczba węzłów rozwiniętych przez A*; miara niezależna
  od implementacji i~sprzętu & A*, Unet-A* \\
  \midrule
Redukcja węzłów & $1 - N_{\mathrm{U}} / N_{\mathrm{A}}$, gdzie
  $N_{\mathrm{U}}$ i~$N_{\mathrm{A}}$ to liczby węzłów rozwiniętych przez
  Unet-A* i~A*; liczona dla połączeń, w~których użyto sieci & Unet-A* \\
\midrule
\multicolumn{3}{l}{\emph{Jakość predykcji}} \\
\midrule
Dice, IoU & zgodność maski przewidzianej przez sieć (próg 0{,}5)
  z~etykietą wyznaczoną algorytmem Lee & Unet-A* \\
\bottomrule
\end{tabular}
\end{table}

Wyniki dla każdej płytki testowej zapisywane są w~osobnym pliku, a~następnie
agregowane dla całego zbioru testowego.

\section{Ocena kompatybilności elektromagnetycznej}
\label{sec:emc}

% EMI jeśli się uda
Środowisko EasyEDA nie udostępnia symulacji pola elektromagnetycznego płytki.
Bezpośrednia ocena emisji zakłóceń wymagałaby zewnętrznego symulatora pola
albo pomiarów fizycznych wykonanych płytek. Z~tego powodu w~badaniach
zastosowano miary pośrednie, wyznaczane z~geometrii ścieżek. Są one zgodne
z~powszechnie stosowanymi zaleceniami projektowania pod kątem EMC:
\begin{itemize}
  \item łączna długość ścieżek --- dłuższe ścieżki skuteczniej promieniują
    i~odbierają zakłócenia;
  \item liczba przelotek --- każda przelotka stanowi nieciągłość impedancji;
  \item długość równoległego prowadzenia ścieżek różnych sieci w~odległości
    mniejszej niż zadany próg --- miara podatności na przesłuch;
  \item pole pętli tworzonej przez ścieżkę sygnałową i~ścieżkę powrotną
    masy --- miara emisji promieniowanej.
\end{itemize}
% openEMS; uzupełnić bibliografię o źródło zaleceń EMC.

\section{Modyfikacje kodu źródłowego}
\label{sec:modyfikacje}

Kod repozytorium Unet-Astar w~udostępnionej postaci nie umożliwiał
przeprowadzenia badań bez zmian. Wprowadzono następujące modyfikacje:
\begin{itemize}
  \item poprawiono wyznaczanie punktu początkowego i~końcowego próbki uczącej
    dla płytek wczytanych z~pliku DSN --- wykorzystywano współrzędne
    nieprzesunięte do początku układu, co dla płytek z~EasyEDA dawało
    ujemne indeksy;
  \item ujednolicono wartość zwracaną przez algorytm A* w~przypadku
    braku ścieżki --- wcześniej powodowała ona przerwanie przetwarzania
    całej płytki;
  \item poprawiono błąd indeksowania sieci w~module wizualizacji, który
    uniemożliwiał trasowanie płytek z~sieciami wielopinowymi;
  \item dodano zapis obiektu płytki wczytanej z~pliku DSN;
  \item dodano skrypty: generowania próbek z~losowych płytek i~z~plików DSN,
    treningu oraz testowania na płytkach z~EasyEDA;
  \item dodano zapis wyników trasowania do plików graficznych i~JSON;
  \item dodano eksport wyniku trasowania do pliku Specctra SES
    importowanego przez EasyEDA.
\end{itemize}

\section{Ograniczenia}
\label{sec:ograniczenia}

\begin{itemize}
  \item Trasowanie odbywa się na siatce 8~mil, a~reguły projektowe są
    zaokrąglane w~górę do jej wielokrotności. Skutkuje to większymi
    odstępami niż wymagane i~może uniemożliwić poprowadzenie ścieżek
    na gęstych płytkach.
  \item Program stosuje jedną szerokość ścieżki i~jeden odstęp dla całej
    płytki --- klasy sieci nie są obsługiwane (tabela~\ref{tab:reguly-ai}).
  \item Ze względu na pojemność pamięci karty graficznej nie wytrenowano
    modelu \emph{ResNetUNet} zaproponowanego przez autorów rozwiązania.
  \item Czasy działania zmierzone dla implementacji w~języku Python nie są
    porównywalne z~czasem działania autoroutera EasyEDA. Dlatego jako
    podstawową miarę wydajności przyjęto liczbę rozwiniętych węzłów.
\end{itemize}
